%%%%%%%%%%  scan idx 67  |  folio 54  %%%%%%%%%%
culier des relations de transitivité auxquelles nous avions fait allusion au \S 2, et qui résultent des suites spectrales des foncteurs composés.

Notons la formule suivante, valable pour $Y$ non singulière (et particulièrement aisée à vérifier lorsque $Y$ est de plus quasi-projective, en prenant pour $y$ la classe d'un fibré vectoriel sur $Y$)~:
\begin{equation}\tag{2.27}
 f_!(x.f^!(y))=f_!(x).y\qquad(x.\in\underline K(X),\ y\in\underline K(Y)).
\end{equation}
Dans le cas particulier où les images réciproques des points de $Y$ par $f$ sont finies, et contenues dans un ouvert affine\footnote{Cette dernière condition est évidemment conséquence de la première (chose que l'auteur ignorait encore en 1957~!).}, on constate que le foncteur $F\longrightarrow f_!(F)$ est exact et que $R^qf_!(F)=0$ pour $q>0$~; dans ce cas, on a donc $f_!(\gamma_X(F))=\gamma_Y(f_!(F))$. En particulier, si $X$ est un sous-espace algébrique de $Y$ et si $i$ est l'injection canonique de $Y$ dans $X$, on a~:
\begin{equation}\tag{2.28}
 i_!(\gamma_X(F))=\gamma_Y(F),
\end{equation}
où, dans le deuxième membre, $F$ est considéré comme un faisceau sur $Y$ nul en dehors de $X$.

\subsection*{\S 4. Quelques résultats techniques}

Soit $X$ un espace algébrique~; on a vu au \S 3 que la projection $p$ de $X\times k$ sur $X$ définit un homomorphisme
\begin{equation}\tag{2.29}
 p^!:\underline K(X)\longrightarrow\underline K(X\times k).
\end{equation}

\textbf{Proposition 2.9.} \emph{L'homomorphisme (2.29) est surjectif.}

On raisonne par récurrence sur $n=\dim X$~; l'assertion est triviale si $n<0$~; supposons $n\geq0$~; en vertu de la proposition 2.6, il suffit de
